\subsection{Research directions in AOSE}
\label{sec:research-directions}
\slidesrc{EMAS keynote: 16--56}

\subsubsection{1. Agent-oriented architectures}
\slidesrc{EMAS keynote: 16--22}

The three families of Section~\ref{sec:agent-architectures}, with their founding papers.
The slides illustrate them with SpongeBob pictures; here are the ideas.
\begin{itemize}
  \item \textbf{Belief-Desire-Intention} (Rao and Georgeff, ``BDI agents: from theory to
        practice'', ICMAS 1995; plus Georgeff and Lansky 1986, Rao and Georgeff 1992,
        Shoham 1993). Slide 17 shows the BDI components: \emph{beliefs} (what the agent
        knows about the environment), \emph{goals}, \emph{plans} (recipes), a
        \emph{reasoner} that selects plans and turns them into \emph{intentions}, which act
        on the environment. Studied in depth in Section~\ref{sec:bdi-agentspeak}.
  \item \textbf{Subsumption} (Brooks, ``A robust layered control system for a mobile
        robot'', 1986). Slide 18: behaviours \emph{Avoid}, \emph{Wander}, \emph{Explore}
        stacked on top of each other, all connected to the environment
        (Section~\ref{sec:reactive-architectures}).
  \item \textbf{Hybrid} (Ferguson's TouringMachines, 1992; M\"uller and Pischel's InteRRaP,
        1993). Slides 19--21 picture a hybrid architecture as a layered cake: several
        horizontal layers, each with its own job, stacked into one agent
        (Section~\ref{sec:deliberative-hybrid}).
\end{itemize}

\begin{nbfact}{FIPA abstract agent platform (slide 22)}
\textbf{FIPA} (Foundation for Intelligent Physical Agents) was created in Switzerland in 1996
and accepted as an IEEE standards committee in 2005. It standardises how agents are hosted
and communicate. A FIPA agent platform contains:
\begin{center}
\begin{tikzpicture}[font=\small, >=Stealth,
  m/.style={draw, thick, rounded corners, minimum height=0.8cm, align=center, fill=#1}]
  \node[m=white] (sw) at (0,3.0) {Software};
  \draw[thick, fill=nbblue!5, rounded corners] (-3.6,0.3) rectangle (3.6,2.3);
  \node[font=\scriptsize\bfseries] at (-2.6,2.05) {Agent platform};
  \node[m=nbgreen!12] (ag) at (-2.3,1.35) {Agent};
  \node[m=nborange!12] (ams) at (0,1.35) {Agent\\Management\\System};
  \node[m=nbpurple!10] (df) at (2.3,1.35) {Directory\\Facilitator};
  \node[m=nbgray!15, minimum width=6.6cm] (mts) at (0,-0.2) {Message Transport System};
  \draw[thick, fill=nbblue!5, rounded corners] (-3.6,-1.9) rectangle (3.6,-1.0);
  \node[m=nbgray!15, minimum width=6.6cm, minimum height=0.6cm] (mts2) at (0,-1.45) {Message Transport System (another platform)};
  \draw[<->, thick] (sw) -- (0,2.3);
  \draw[<->, thick] (ag) -- (ag |- mts.north);
  \draw[<->, thick] (ams) -- (ams |- mts.north);
  \draw[<->, thick] (df) -- (df |- mts.north);
  \draw[<->, thick] (mts) -- (mts2);
\end{tikzpicture}
\end{center}
\begin{itemize}[nosep]
  \item \textbf{Agent Management System} (AMS): the ``white pages'': creates and destroys
        agents, knows every agent's name and address;
  \item \textbf{Directory Facilitator} (DF): the ``yellow pages'': agents register the
        \emph{services} they offer and search for providers;
  \item \textbf{Message Transport System}: delivers messages, also to agents on other
        platforms.
\end{itemize}
JADE (Section~\ref{sec:jade}) is a FIPA-compliant platform with exactly these components.
\end{nbfact}

\subsubsection{2. AOSE methodologies and notation techniques}
\slidesrc{EMAS keynote: 23--29}

Many agent-oriented methodologies exist (Henderson-Sellers and Giorgini, eds.,
\emph{Agent-Oriented Methodologies}, 2005). Slide 24 shows their genealogy: most derive
from object-oriented methods (OO, RUP, OMT, Fusion, OPEN) or from knowledge engineering
(MAS-CommonKADS); among them Gaia, Prometheus, Tropos, MaSE, PASSI, INGENIAS, MESSAGE,
Agent OPEN, ADELFE.

\begin{nbdef}{Gaia (Wooldridge, Jennings, Kinny, 2000)}
A methodology for agent-oriented \textbf{analysis and design}. It sees a MAS as an
\emph{organisation} of interacting \emph{roles}, and proceeds through phases:
\begin{center}
\begin{tikzpicture}[font=\scriptsize, >=Stealth, node distance=0.4cm,
  m/.style={draw, thick, rounded corners, minimum width=2.1cm, minimum height=0.65cm, align=center, fill=#1}]
  \node[draw, ellipse, thick] (req) at (0,0) {Requirements};
  \node[m=nbblue!10] (env) at (0,-1.1) {Environmental model};
  \node[m=nbblue!10] (prm) at (-2.7,-2.0) {Preliminary\\role model};
  \node[m=nbblue!10] (pim) at (2.7,-2.0) {Preliminary\\interaction model};
  \node[m=nbblue!10] (or) at (0,-2.6) {Organisational\\rules};
  \node[m=nbpurple!10] (os) at (0,-3.8) {Organisational\\structure};
  \node[m=nbpurple!10] (rm) at (-2.7,-4.8) {Role model};
  \node[m=nbpurple!10] (im) at (2.7,-4.8) {Interaction model};
  \node[m=nborange!12] (am) at (-2.7,-6.1) {Agent model};
  \node[m=nborange!12] (sm) at (2.7,-6.1) {Services model};
  \draw[->, thick] (req) -- (env);
  \draw[->, thick] (env) -- (prm); \draw[->, thick] (env) -- (pim);
  \draw[->, thick] (prm) -- (or); \draw[->, thick] (pim) -- (or);
  \draw[->, thick] (or) -- (os);
  \draw[->, thick] (os) -- (rm); \draw[->, thick] (os) -- (im);
  \draw[->, thick] (prm) -- (rm); \draw[->, thick] (pim) -- (im);
  \draw[->, thick] (rm) -- (am); \draw[->, thick] (rm) -- (sm);
  \draw[->, thick] (im) -- (sm); \draw[->, thick] (im) -- (am);
  \draw[->, thick] (am) -- ++(0,-0.8); \draw[->, thick] (sm) -- ++(0,-0.8);
  \foreach \y/\t in {-0.5/{\textbf{Analysis}}, -3.25/{\textbf{Architectural design}}, -5.5/{\textbf{Detailed design}}, -6.95/{\textbf{Implementation (free)}}}{
    \draw[dashed, nbgray] (-4.4,\y) -- (4.4,\y); \node[anchor=west, nbgray!80!black] at (4.5,\y-0.3) {\t};
  }
\end{tikzpicture}
\end{center}
Gaia deliberately leaves the \textbf{implementation} stage free: the design can be
implemented with an agent platform or with mainstream technologies (the point made in
Section~\ref{sec:uses-of-mas}).
\end{nbdef}

\paragraph{Notations.} Agents need diagrams that UML does not offer. \textbf{AUML} (Agent
UML; Odell, Parunak, Bauer, \emph{Extending UML for agents}, 2000) adds, among other things,
\textbf{protocol diagrams}; \textbf{Prometheus} (Winikoff and Padgham, 2004) has its own
icons for agents, actions, messages, data, percepts, protocols, capabilities and plans
(slide 29).

\begin{nbfact}{The FIPA Request interaction protocol (slides 26--27: ``Remember this protocol!'')}
\begin{center}
\begin{tikzpicture}[font=\small, >=Stealth, x=1cm, y=-0.72cm]
  \node[draw, thick, fill=nbblue!8, minimum width=2.4cm] (I) at (0,0) {Initiator};
  \node[draw, thick, fill=nbgreen!8, minimum width=2.4cm] (P) at (7,0) {Participant};
  \draw[dashed, nbgray] (0,0.6) -- (0,8.4); \draw[dashed, nbgray] (7,0.6) -- (7,8.4);
  \draw[->, thick] (0,1.3) -- node[above]{\code{request}} (7,1.3);
  \draw[->, thick, nbred] (7,2.3) -- node[above]{\code{refuse}} (0,2.3);
  \node[nbred, font=\scriptsize] at (3.5,2.75) {[refused]};
  \draw[->, thick, nbgreen!60!black] (7,4.0) -- node[above]{\code{agree}} (0,4.0);
  \node[nbgreen!60!black, font=\scriptsize] at (3.5,4.45) {[agreed and notification necessary]};
  \draw[->, thick, nbred] (7,5.5) -- node[above]{\code{failure}} (0,5.5);
  \draw[->, thick, nbblue] (7,6.6) -- node[above]{\code{inform-done : inform}} (0,6.6);
  \draw[->, thick, nbblue] (7,7.7) -- node[above]{\code{inform-result : inform}} (0,7.7);
  \draw[decorate, decoration={brace, amplitude=5pt, mirror}, thick, nbgray] (7.3,2.1) -- node[right=5pt, font=\scriptsize, align=left]{one of\\these two} (7.3,4.0);
  \draw[decorate, decoration={brace, amplitude=5pt, mirror}, thick, nbgray] (7.3,5.3) -- node[right=5pt, font=\scriptsize, align=left]{if agreed:\\one of\\these three} (7.3,7.9);
\end{tikzpicture}
\end{center}
\begin{enumerate}[nosep]
  \item The initiator sends a \code{request} (asks the participant to perform an action).
  \item The participant answers \code{refuse} (protocol ends) or \code{agree}.
  \item If it agreed, it performs the action and finally sends one of: \code{failure} (it
        tried but failed), \code{inform-done} (done, nothing to return) or
        \code{inform-result} (done, here is the result).
\end{enumerate}
(The \code{agree} can be skipped if the participant answers immediately with the final
message.) This protocol is used again in agent communication
(Section~\ref{sec:agent-communication}) and in JADE.
\end{nbfact}

\begin{nbexample}{the Request protocol in a pizzeria}
\notinslides
Customer agent: \code{request}(``deliver one margherita''). Pizzeria agent: \code{agree}
(``ok, 30 minutes''). Later: \code{inform-done} (``delivered''); or \code{failure} (``the
scooter broke down''). If the pizzeria were closed: \code{refuse} (``closed on Mondays'').
\end{nbexample}

\begin{nbfact}{FIPA ACL message structure (slide 28)}
\begin{center}
\small
\begin{tabular}{@{}ll@{}}
\toprule
\textbf{Parameter} & \textbf{Category} \\
\midrule
performative & type of communicative act \\
sender, receiver, reply-to & participant in communication \\
content & content of message \\
language, encoding & description of content \\
\bottomrule
\end{tabular}

\smallskip
\begin{tabularx}{\linewidth}{@{}l X@{}}
\toprule
\textbf{Performative} & \textbf{Meaning} \\
\midrule
\code{REQUEST} & the sender asks the receiver to perform some action \\
\code{AGREE} & the sender agrees to perform the requested action \\
\code{REFUSE} & the sender refuses to perform the requested action \\
\code{NOT\_UNDERSTOOD} & someone participating in the conversation didn't understand something \\
\code{INFORM} & a positive reply because an action is performed, or the answer to some previous question \\
\code{FAILURE} & a failure in performing the requested action happened after an initial agreement \\
\bottomrule
\end{tabularx}
\end{center}
Studied in detail in Section~\ref{sec:agent-communication}.
\end{nbfact}

\subsubsection{3. Agent-oriented programming languages}
\slidesrc{EMAS keynote: 30--39}

\paragraph{Logic-based agents (slides 31--32).} The ``traditional'' approach to building
artificial intelligent systems: \textbf{logical formulae} give a symbolic representation of
the agent's environment and desired behaviour, and \textbf{logical deduction} (theorem
proving) syntactically manipulates this representation. Slide 31's cartoon thinks the Prolog
program \code{mortal(X) :- human(X).} \code{human(socrates).} with the query
\code{?- mortal(socrates).} answered \code{yes} (Whiteboard~\ref{wb:modus-ponens}).

\paragraph{AgentSpeak(L) (slides 33--36).} The BDI architecture is one of the best known and
most studied architectures for \emph{cognitive} agents; \textbf{AgentSpeak(L)} (A.~S.~Rao,
``AgentSpeak(L): BDI agents speak out in a logical computable language'', MAAMAW 1996) is a
logic-based programming language inspired by it:
\begin{itemize}[nosep]
  \item \textbf{beliefs} represent the knowledge of the agent;
  \item \textbf{goals} are beliefs, or conjunctions and disjunctions of beliefs, which must
        be \emph{achieved} or \emph{tested} in the current state;
  \item \textbf{plans} contain the \emph{procedural} knowledge of the agent, characterised by
        a trigger, a context, a body, a maintenance condition, a set of ``success actions''
        and a set of ``failure actions'';
  \item \textbf{intentions} are partially instantiated plans.
\end{itemize}
Its syntax (slide 35) and semantics (slide 36) are explained in
Section~\ref{sec:bdi-agentspeak}.

\paragraph{JadeScript (slides 37--39).} The agent-oriented programming language for JADE
agents (Bergenti, Caire, Monica, Poggi, 2020). A Ping--Pong example: \code{Ping} sends a
\code{request} for the action \code{Reply}; \code{Pong} answers with an \code{inform} that
the proposition \code{Received} holds.
\begin{minted}[linenos=false]{text}
ontology PingPong
    action Reply
    proposition Received

agent Pong
    uses ontology PingPong
    on create do
        log "Agent 'Pong' created."
        activate WaitAndReply

cyclic behavior WaitAndReply for agent Pong
    on message request Reply do
        sender = sender of message
        log "received 'Reply' request from agent '" + name of sender + "'"
        inform Received to sender

agent Ping
    uses ontology PingPong
    on create do
        log "Agent 'Ping' created."
        activate WaitFromPong
        activate SendRequest

one shot behavior SendRequest for agent Ping
    on execute do
        log "sending 'Reply' request..."
        send message request Reply to "Pong"@
\end{minted}
Notice the agent concepts as first-class language keywords: \emph{agent},
\emph{ontology} (the shared vocabulary), \emph{behavior} (cyclic or one-shot, as in JADE,
Section~\ref{sec:jade}), \emph{message} with a performative (\code{request},
\code{inform}).

\subsubsection{4. MAS infrastructures and toolkits}
\slidesrc{EMAS keynote: 40--46}

Slide 41 lists the papers of the \emph{Agent Toolkits} special track at EUMAS 2024 (VEsNA,
belief revision in JavaScript, protocol-based MAS, Yggdrasil for hypermedia MAS, SARL,
Jadescript, JaKta, Jason+BSPL, JaCaMo, Jadex V, learning and explaining in BDI agents): the
field is alive. Two infrastructures are highlighted:
\begin{itemize}
  \item \textbf{JaCaMo} (2007--2020+): a platform combining \textbf{Ja}son (agents),
        \textbf{Ca}rtAgO (artifacts: the environment) and \textbf{Mo}ise (organisations).
  \item \textbf{JADE} (book \emph{Developing multi-agent systems with JADE}, 2007) ---
        ``Remember this framework!'' (slide 45). Section~\ref{sec:jade}.
\end{itemize}

\subsubsection{5. Agent-based modelling and simulation}
\slidesrc{EMAS keynote: 47--48}

Pointer to a summer-school tutorial on running agent-based models (e.g.\ NetLogo) on
high-performance computing facilities. See Section~\ref{sec:abms}.

\subsubsection{6. Integration of symbolic and sub-symbolic approaches}
\slidesrc{EMAS keynote: 49--56}

The keynote repeats the System~1 / System~2 argument of Section~\ref{sec:symbolic-vs-subsymbolic}
(Kahneman 2011; Booch et al., \emph{Thinking Fast and Slow in AI}, 2021) and shows that the
integration of symbolic and sub-symbolic (``neuro-symbolic'') approaches is a large research
area (Google Scholar searches, slides 54--55). It is not new: slide 56 shows J.~Shavlik,
\emph{Combining Symbolic and Neural Learning} (Machine Learning, 1994), already reviewing
hybrids of the two approaches thirty years ago. For agents, the typical combination is a BDI
agent (symbolic reasoning, explainable plans) using learned components for perception,
language understanding or plan selection.
